\documentclass[12pt,a4paper]{article}
\usepackage[utf8]{inputenc}
\usepackage{amsmath,amssymb}
\usepackage{graphicx}
\usepackage[left=1.5cm,right=1cm,top=1.3cm,bottom=1.3cm]{geometry}
\usepackage{lastpage}
% Font Times New Roman (cần XeLaTeX + fontspec — uncomment nếu VPS hỗ trợ):
% \usepackage{fontspec}
% \setmainfont{Times New Roman}

\begin{document}

\begin{center}
\begin{tabular}{p{0.4\textwidth}p{0.6\textwidth}}
\textbf{1} & 5 \tabularnewline
\textbf{2} & \textbf{6} \tabularnewline
\textbf{3} & \textit{7, 8} \tabularnewline
\textit{4} & \tabularnewline
Họ và tên: ...................................................................................... \newline Số báo danh: .................................................................................. & \textbf{Mã đề: 1761} \tabularnewline
\end{tabular}
\end{center}
\vspace{0.5cm}

\noindent\textbf{PHẦN I. Câu trắc nghiệm nhiều phương án lựa chọn.}
Thí sinh trả lời từ câu 1 đến câu 2. Mỗi câu thí sinh chỉ chọn một phương án.

\medskip

\noindent\textcolor{blue}{\textbf{Câu 1.}} Cho hàm số $f(x) = \dfrac{1}{x} + \dfrac{1}{\cos ^2 x}$. Chọn khẳng định đúng.

\noindent\begin{tabular}{p{\linewidth}}
\textbf{A.}~$\displaystyle\int f(x) \mathrm{\,d}x = \ln |x| + \tan x + C$. \\
\textbf{B.}~$\displaystyle\int f(x) \mathrm{\,d}x = \ln x + \tan x + C$. \\
\textbf{C.}~$\displaystyle\int f(x) \mathrm{\,d}x = \ln x + \tan |x| + C$. \\
\textbf{D.}~$\displaystyle\int f(x) \mathrm{\,d}x = \ln |x| + \tan x$. \\
\end{tabular}

\noindent\textcolor{blue}{\textbf{Câu 2.}} 
% [Hình TikZ không compile được: tikz_e0c4e8580c64]
Cho hàm số $y=f(x)$ có đồ thị như hình vẽ. Hàm số đồng biến trên khoảng nào đưới đây?

\noindent\begin{tabular}{p{0.25\linewidth} p{0.25\linewidth} p{0.25\linewidth} p{0.25\linewidth}}
\textbf{A.}~$(0;+\infty)$. & \textbf{B.}~$(-2;2)$. & \textbf{C.}~$(-1;1)$. & \textbf{D.}~$(-2;1)$. \\
\end{tabular}

\noindent\textbf{PHẦN II. Câu trắc nghiệm đúng sai.}
Thí sinh trả lời từ câu 1 đến câu 2. Trong mỗi ý, thí sinh chọn đúng (Đ) hoặc sai (S).

\medskip

\noindent\textcolor{blue}{\textbf{Câu 1.}} 
% [Hình TikZ không compile được: tikz_2c10d18afdee]
Cho hàm số $y=f(x)=\dfrac{x^2+bx+c}{x-2}$ có đạo hàm $f'(x)$. Đồ thị của hàm số $f'(x)$ như hình vẽ bên.

\noindent a)~Phương trình $f'(x)=0$ có hai nghiệm $x=1$ và $x=3$.

\noindent b)~Hàm số $y=f(x)$ nghịch biến trên khoảng $(1;3)$.

\noindent c)~Hàm số $y=f(x)$ đạt cực đại tại $x=1$ và đạt cực tiểu tại $x=3$.

\noindent d)~Nếu $f(0)=1$ thì $\max\limits_{[3;4]} f(x)=6$.

\noindent\textcolor{blue}{\textbf{Câu 2.}} Một vật chuyển động thẳng được cho bởi phương trình $s(t)=-\dfrac{1}{3}t^3+4t^2+9t$, trong đó $t$ tính bằng giây và $s(t)$ là vị trí của vật tại thời điểm $t$, tính bằng mét.

\noindent a)~Vận tốc của vật tại các thời điểm $t=3$ giây là $v(3)=1 \text{ m/s}$.

\noindent b)~Quãng đường vật đi được từ lúc bắt đầu chuyển động đến khi vật đứng yên là $162 \text{ m}$.

\noindent c)~Gia tốc của vật tại thời điểm $t=3$ giây là $a(3)=2 \text{ m/s}^2$.

\noindent d)~Vật tăng tốc khi $t \in[0; 4]$.

\noindent\textbf{PHẦN III. Câu trả lời ngắn.}
Thí sinh trả lời từ câu 1 đến câu 2.

\medskip

\noindent\textcolor{blue}{\textbf{Câu 1.}} Cho hàm số $y=\sqrt{1-x^2}$. Gọi $(a;b)$ là khoảng đồng biến có độ dài lớn nhất của hàm số trên. Tính $a+b$.

\noindent\textcolor{blue}{\textbf{Câu 2.}} Hàm số $y=\log_2(x^2+3x)$ nghịch biến trên khoảng $(-\infty; a)$. Giá trị lớn nhất của $a$ là bao nhiêu?
	\par

\noindent\textbf{PHẦN IV. Câu tự luận.}
Thí sinh trả lời từ câu 1 đến câu 2.

\medskip

\noindent\textcolor{blue}{\textbf{Câu 1.}} [Trích đề ôn tập HKI-Lớp 12-THPT Bùi Thị Xuân-Năm học 2024-2025]
	Một vật đang đứng yên thì bắt đầu chuyển động theo quy luật $s(t)=-t^3+3t^2+6t$, với $t$ (giây) là khoảng thời gian tính từ lúc vật bắt đầu chuyển động và $s$ (mét) là quãng đường vật đi được trong khoảng thời gian đó. Hỏi vật tăng tốc trong khoảng thời gian bao nhiêu giây tính từ lúc bắt đầu chuyển động?

\noindent\textcolor{blue}{\textbf{Câu 2.}} Sau khi phát hiện một bệnh dịch, các chuyên gia y tế ước tính số người nhiễm bệnh kể từ ngày xuất hiện bệnh nhân đầu tiên đến ngày thứ $t$ là $f(t)=45t^2-t^3$, $t=0$, $1$, $2$, $\ldots$, $25$. Nếu coi $f(t)$ là hàm số xác định trên đoạn $[0; 25]$ thì đạo hàm $f'(t)$ được xem là tốc độ truyền bệnh (người/ngày) tại thời điểm $t$. Giả sử khoảng thời gian mà tốc độ truyền bệnh giảm là từ ngày thứ $m$ đến ngày thứ $n$. Khi đó $n-m$ bằng bao nhiêu?


\textbf{------------ HẾT ------------}
\label{lastpage}

\end{document}